\documentclass[11pt]{article}
\usepackage{amsmath,amssymb,amsthm,mathrsfs}
\usepackage[margin=1.1in]{geometry}
\usepackage{hyperref}

\theoremstyle{plain}
\newtheorem{theorem}{Theorem}
\newtheorem{lemma}[theorem]{Lemma}
\newtheorem{proposition}[theorem]{Proposition}
\newtheorem{corollary}[theorem]{Corollary}
\theoremstyle{definition}
\newtheorem{definition}[theorem]{Definition}
\newtheorem{hypothesis}[theorem]{Hypothesis}
\theoremstyle{remark}
\newtheorem{remark}[theorem]{Remark}

\newcommand{\R}{\mathbb{R}}
\newcommand{\C}{\mathbb{C}}
\newcommand{\Z}{\mathbb{Z}}
\newcommand{\Lam}{\Lambda}
\DeclareMathOperator{\spec}{spec}
\DeclareMathOperator{\rank}{rank}
\DeclareMathOperator{\inertia}{inertia}
\DeclareMathOperator{\ran}{ran}

\title{The Gram double pencil: a spectral census for the zeros of $\zeta$,\\
and an Euler-anchored regulator}
\author{}
\date{2026-08-14}

\begin{document}
\maketitle

\begin{abstract}
We attach to each window of the critical line a pair of Hankel matrices computed by
contour integration of $\xi'/\xi$, and prove that the number of zeros off the critical
line in the window equals the negative inertia of the first matrix (Theorem
\ref{thm:inertia}); the remaining spectral data of the pair counts the zeros with
multiplicity, their distinct support, and their ordinates (Theorem \ref{thm:census}).
Both statements are unconditional. The Riemann Hypothesis is thereby equivalent to a
Stieltjes-type positivity condition, window by window, in a quotient chart obtained
from the functional equation (Proposition \ref{prop:stieltjes}). We then construct,
in the half plane of absolute convergence and with no reference to zeros, an
Euler-anchored double pencil which is unconditionally positive semidefinite
(Proposition \ref{prop:anchor}), and prove that its spectrum can be transported
isospectrally by an explicit finite-dimensional regulator whenever a bank response
matrix is nonsingular (Theorem \ref{thm:regulator}), extended through spectral
collisions by a block regulator (Lemma \ref{lem:block}). We then construct the readout
itself: the warp --- the flow of a bank-computable velocity field --- exists, is real
and order-preserving, and carries the frozen labels exactly onto the moving spectrum
(Theorem \ref{thm:warp}); the characteristic polynomial obeys an exact
determinant-pullback law with nonvanishing cofactor (Theorem \ref{thm:pullback}). The
program thereby reduces to a single hypothesis, the seat (Hypothesis \ref{hyp:seat}):
terminal seating of the warped anchor spectrum on the window support, from which the
positivity, the census and the critical line follow (Theorem \ref{thm:chain}). The
remaining obligations are stated precisely in \S\ref{sec:obligations}.
\end{abstract}

\section{Setup}

Let $\xi$ be the completed Riemann zeta function and put
\[
  A(z) := \xi\!\left(\tfrac12 - iz\right),
\]
an entire function of order one, even ($A(-z)=A(z)$, the functional equation), real on
$\R$, whose zeros are $z_\rho = \gamma_\rho + i(\tfrac12-\beta_\rho)\cdot(-1)$ for
$\rho = \beta_\rho + i\gamma_\rho$ a nontrivial zero of $\zeta$; in particular
\[
  \rho \text{ on the critical line} \iff z_\rho \in \R .
\]
Nonreal zeros of $A$ occur in conjugate pairs $z,\bar z$ with equal multiplicities
(conjugate symmetry of $\xi$ on the real axis).

\begin{definition}[window moments]\label{def:moments}
Let $W=(a,b)$ with $A$ nonvanishing on $\partial W \times \{|y|\le Y\}$, $Y>\tfrac12$,
and let $\partial R_W$ be the positively oriented rectangle boundary. Put
\[
  \mu_k(W) := \frac{1}{2\pi i}\oint_{\partial R_W} z^k\,\frac{A'(z)}{A(z)}\,dz,
  \qquad k \ge 0,
\]
and for $n\ge 1$,
\[
  H_n(W) := \bigl(\mu_{\alpha+\beta}(W)\bigr)_{\alpha,\beta=0}^{n-1},
  \qquad
  H_n^{(1)}(W) := \bigl(\mu_{\alpha+\beta+1}(W)\bigr)_{\alpha,\beta=0}^{n-1}.
\]
\end{definition}

By the argument principle, with $x_1,\dots,x_r$ the distinct real zeros of $A$ in the
window (multiplicities $a_i\ge 1$) and $z_j,\bar z_j$, $j=1,\dots,q$, the distinct
conjugate pairs (equal multiplicities $b_j \ge 1$),
\begin{equation}\label{eq:moments}
  \mu_k(W) \;=\; \sum_{i=1}^r a_i x_i^k \;+\; \sum_{j=1}^q b_j\bigl(z_j^k + \bar z_j^{\,k}\bigr).
\end{equation}
All $\mu_k(W)$ are real and computable from $A$ alone; no knowledge of the zeros is used.
Write $m := r + 2q$ for the number of distinct support points.

\section{The inertia identity}

\begin{theorem}[inertia]\label{thm:inertia}
For every $n \ge m$,
\[
  \inertia H_n(W) \;=\; \bigl(r+q,\; q,\; n-r-2q\bigr),
\]
where $\inertia$ denotes (positive, negative, zero) eigenvalue counts. In particular
\[
  n_-\bigl(H_n(W)\bigr) \;=\; q,
  \qquad\text{and}\qquad
  H_n(W) \succeq 0 \iff q = 0,
\]
i.e.\ the negative inertia counts exactly the off-critical-line conjugate pairs.
\end{theorem}

\begin{proof}
Let $v(t) := (1,t,\dots,t^{n-1})^{T}$. By \eqref{eq:moments},
\[
  H_n \;=\; \sum_{i=1}^r a_i\, v(x_i)v(x_i)^{T}
  \;+\; \sum_{j=1}^q b_j\Bigl[v(z_j)v(z_j)^{T} + v(\bar z_j)v(\bar z_j)^{T}\Bigr].
\]
Write $v(z_j) = A_j + iB_j$ with $A_j,B_j\in\R^n$; then $v(\bar z_j)=A_j-iB_j$ and
\begin{equation}\label{eq:minus}
  v(z_j)v(z_j)^{T} + v(\bar z_j)v(\bar z_j)^{T}
  \;=\; 2A_jA_j^{T} - 2B_jB_j^{T},
\end{equation}
the minus sign being the whole mechanism. Hence $H_n = UJU^{T}$ with
\[
  U = \bigl[v(x_1)\cdots v(x_r)\;A_1\,B_1\cdots A_q\,B_q\bigr],
  \qquad
  J = \operatorname{diag}(a_1,\dots,a_r,\,2b_1,-2b_1,\dots,2b_q,-2b_q),
\]
so $\inertia J = (r+q,q,0)$.

$U$ has full column rank $m$: the complex matrix
$V=[v(x_1),\dots,v(x_r),v(z_1),v(\bar z_1),\dots]$ is Vandermonde on $m$ distinct nodes
with $n\ge m$, so its leading $m\times m$ minor equals
$\prod_{\alpha<\beta}(s_\beta-s_\alpha)\ne 0$ and $\rank V=m$; replacing each pair
$(v(z_j),v(\bar z_j))$ by $(A_j,B_j) = \bigl(\tfrac{v(z_j)+v(\bar z_j)}{2},
\tfrac{v(z_j)-v(\bar z_j)}{2i}\bigr)$ is an invertible change of columns, so
$\rank U = m$.

Take a thin $QR$ factorisation $U=QR$ with $Q^{T}Q=I_m$ and $R$ invertible, and extend
$Q$ to an orthogonal $\widetilde Q = [Q,Q_\perp]$. Then
\[
  \widetilde Q^{T} H_n \widetilde Q =
  \begin{pmatrix} RJR^{T} & 0\\ 0 & 0_{n-m}\end{pmatrix},
\]
and by Sylvester's law of inertia $\inertia(RJR^{T}) = \inertia(J)$. The zero block
contributes $n-m$ null directions.
\end{proof}

\begin{remark}
The hypothesis $n \ge m$ is necessary: an off-line pair is invisible to a Hankel matrix
smaller than the distinct support count. Since $\mu_0(W)$ is the number of zeros in $W$
with multiplicity and $m \le \mu_0(W)$, the choice
$n := \lceil \mu_0(W)\rceil + 1$ is admissible and computable from the bank alone.
\end{remark}

\begin{remark}[collisions]
If a conjugate pair approaches $\R$, the corresponding negative eigenvalue tends to $0^-$
and at collision becomes a null direction; the invariant through such an event is the
inertia, not any individually ordered eigenvalue. Consequently the correct target is
non-negativity $H_n \succeq 0$ (equivalently $n_-=0$), never strict positivity: the
$n-m$ structural nulls and the collision nulls are expected.
\end{remark}

\section{The census, and the equivalent form of RH}

\begin{theorem}[spectral census]\label{thm:census}
Let $n \ge m$. From the spectral data of $H_n(W)$ and the value $\mu_0(W)$:
\[
  \mu_0 = \text{number of zeros in } W \text{ with multiplicity},\qquad
  \rank H_n = m = r+2q,
\]
\[
  n_-(H_n) = q,\qquad n_+(H_n)-n_-(H_n) = r,
\]
and $\mu_0 > \sum_i 1 + \sum_j 2$ (i.e.\ $\mu_0$ exceeds the distinct count) if and only
if some zero in $W$ is multiple. Moreover the generalized eigenvalues of the pencil
$\bigl(H^{(1)}_n, H_n\bigr)$ restricted to $\ran H_n$ are exactly the support points
$x_i$ and $z_j,\bar z_j$.
\end{theorem}

\begin{proof}
The first four items are Theorem \ref{thm:inertia} together with $\mu_0=\sum a_i+2\sum b_j$.
For the pencil statement, $H_n = UJU^T$ and $H_n^{(1)} = U J \Sigma U^{T}$ where $\Sigma$
is the diagonal matrix of support points in the basis of columns of $U$ (Vandermonde
shift identity $v(t)^{T}\!\cdot t = $ shifted moments); on $\ran H_n$ the pencil is
therefore similar to $\Sigma$. This is Gaussian quadrature for the measure
\eqref{eq:moments}.
\end{proof}

\begin{proposition}[quotient chart; Stieltjes form]\label{prop:stieltjes}
Let $w=z^2$ and $\nu_k(W) := \mu_{2k}(W)$. Under $z\mapsto w$: a real zero maps to
$w>0$; a conjugate pair $x\pm iy$ ($x\ne0$) maps to a conjugate pair in $w$; a pair on
the imaginary axis maps to $w<0$. Consequently, for $n$ large enough in the sense of
Theorem \ref{thm:inertia},
\[
  \text{all zeros of } A \text{ in } W \text{ are real}
  \iff
  \bigl(\nu_{\alpha+\beta}\bigr) \succeq 0
  \ \text{ and }\
  \bigl(\nu_{\alpha+\beta+1}\bigr) \succeq 0 ,
\]
the Stieltjes condition for the measure $\sum_\rho m_\rho \delta_{z_\rho^2}$. Hence RH is
equivalent to: for every window of a tiling of $\R$, both Hankel matrices of the even
moments are positive semidefinite.
\end{proposition}

\begin{proof}
The functional equation makes $A$ even, so the multiset of zeros is symmetric under
$z\mapsto -z$ and descends to the $w$-chart. Reality and non-negativity of the
$w$-support is exactly the Stieltjes moment condition (Hamburger positivity of both
$(\nu_{\alpha+\beta})$ and its shift). The equivalence with the location statement is
Theorem \ref{thm:inertia} applied in the $w$-chart together with the support map.
\end{proof}

\begin{corollary}[sign rigidity]\label{cor:rigidity}
Let $t\mapsto A_t$ be a continuous deformation with zeros varying continuously and no
zero crossing $\partial W$. Then $H_n(t)$ is continuous, and a conjugate pair is created
or annihilated at $t_0$ if and only if an eigenvalue of $H_n(t)$ crosses $0$ at $t_0$;
at simultaneous collisions the invariant is the inertia change. Consequently, if
$H_n(t)\succeq 0$ for all $t$, then $n_-(H_n(t))\equiv 0$ and no off-line pair is ever
born. If moreover $H_n(t_\ast)\succeq0$ is known at one parameter $t_\ast$ (for instance
in a numerically verified range), the conclusion propagates along the deformation.
\end{corollary}

\section{Unconditional statements}

The following hold with no hypothesis.

\begin{enumerate}
\item[(U1)] For every window $W$ and $n\ge m(W)$: $q(W)=n_-\bigl(H_n(W)\bigr)$; the
off-critical-line pair count of a window is the negative inertia of a matrix computed by
contour integration of $A'/A$.
\item[(U2)] The census of Theorem \ref{thm:census} holds as stated.
\item[(U3)] RH is equivalent to the Stieltjes condition of Proposition
\ref{prop:stieltjes} on every window of a tiling, equivalently (Corollary
\ref{cor:rigidity}) to the absence of a negative-sector eigenvalue crossing under the
registered deformation.
\item[(U4)] For every window contained in the numerically verified range
($|t|\le 3\cdot 10^{12}$ by the computations of Platt and Trudgian
\cite{PlattTrudgian2021}), $q(W)=0$ and hence
$H_n(W)\succeq 0$: the positivity hypothesis is an established fact on an initial
segment of the tiling.
\item[(U5)] With $Y>\tfrac12$ the top edge of $\partial R_W$ lies in the half plane of
absolute convergence and the bottom edge maps there by the functional equation, giving a
decomposition $H_n = H^{\mathrm{pr}} + H^{\mathrm{arch}} + H^{\mathrm{side}}$ into blocks
computable from Dirichlet data, $\Gamma$-factors and the lateral edges respectively; in a
tiling, adjacent lateral contributions cancel pairwise.
\end{enumerate}

\section{The Euler-anchored pencil and its regulator}

Nothing in this section refers to zeros; the dependency order is
\[
  \text{Euler/FE bank}\;\to\;(G_1,G_0)\;\to\;\text{real spectrum}\;\to\;
  \text{contour transport}\;\to\;\text{residues}\;\to\;\text{zeros},
\]
and no step may be inverted.

\begin{proposition}[Euler anchor]\label{prop:anchor}
Fix $s_0>1$ and $N\ge1$, and set
\[
  G_\ell[j,k] \;:=\; \sum_{n\ge2} \Lam(n)\,n^{-s_0}\,(\log n)^{\,j+k+\ell},
  \qquad 0\le j,k<N,\quad \ell\in\{0,1\}.
\]
The series converge absolutely, and $G_0,G_1$ are the Hankel matrices of the positive
measure $\sum_n \Lam(n)n^{-s_0}\delta_{\log n}$ supported in $[\log 2,\infty)$. Hence
unconditionally
\[
  G_0 \succeq 0,\qquad G_1 \succeq 0,
\]
and $\spec(G_1,G_0)\subset[\log2,\infty)$ consists of the Gauss quadrature nodes of that
measure.
\end{proposition}

\begin{proof}
Absolute convergence for $s_0>1$ is standard. Positivity is the Hamburger/Stieltjes
criterion for a positive measure on $[\log2,\infty)$: for $c\in\R^N$ and
$P(x)=\sum_j c_jx^j$, $c^{T}G_0c=\int P^2\,d\sigma\ge0$ and
$c^{T}G_1c=\int xP^2\,d\sigma\ge0$.
\end{proof}

\begin{theorem}[unconditional reverb regulator]\label{thm:regulator}
Let $t\mapsto (G_0(t),G_1(t))$ be a smooth family with $G_0(t)>0$ on the active quotient,
and set $S(t):=G_0^{-1/2}G_1G_0^{-1/2}$, so $S=S^{*}$ and
$\spec S=\spec(G_1,G_0)$. Suppose the spectrum of $S(t_0)$ is simple, with spectral
projections $P_a$ and eigenvectors $v_a$, and let the transport be controlled,
\[
  \dot S \;=\; F(u) \;=\; F_0 + \sum_{r=1}^{m} u_r F_r ,
  \qquad F_r = F_r^{*}.
\]
Define $M_{ar} := \langle v_a, F_r v_a\rangle$ and $b_a := -\langle v_a,F_0v_a\rangle$.
If $\det M \ne 0$, then $u := M^{-1}b$ is the unique control with
$\Delta\bigl(F(u)\bigr):=\sum_a P_aF(u)P_a = 0$, and with
\[
  K := \sum_{a\ne b}\frac{P_aF(u)P_b}{\lambda_a-\lambda_b},
  \qquad K^{*}=-K,
\]
one has $\dot S = [S,K]$. Consequently $S(t)=U(t)S(t_0)U(t)^{*}$ with $\dot U=-KU$,
$U$ unitary, and
\[
  \spec S(t) \;=\; \spec S(t_0)\quad\text{exactly.}
\]
In particular the transported pencil remains Hermitian-definite and its generalized
spectrum remains real.
\end{theorem}

\begin{proof}
First-order perturbation theory gives $\dot\lambda_a=\langle v_a,F v_a\rangle$, so
$\Delta(F)$ is precisely the component of the flow able to move eigenvalues; the linear
system $Mu=b$ expresses $\Delta(F(u))=0$ and is uniquely solvable iff $\det M\ne0$.
For such $u$, using $\sum_aP_a=I$,
\[
  [S,K]=\sum_{a\ne b}(\lambda_a-\lambda_b)\frac{P_aF(u)P_b}{\lambda_a-\lambda_b}
  =\sum_{a\ne b}P_aF(u)P_b=F(u)-\Delta(F(u))=F(u)=\dot S .
\]
$K^{*}=-K$ since $F(u)^{*}=F(u)$ and the denominators are real and antisymmetric in
$(a,b)$. The solution of $\dot U=-KU$, $U(t_0)=I$, is unitary, and
$\frac{d}{dt}\bigl(U^{*}SU\bigr)=U^{*}\bigl(\dot S-[S,K]\bigr)U=0$.
\end{proof}

\begin{remark}[numerical verification]
For the Euler anchor of Proposition \ref{prop:anchor} with $s_0=1.5$, $N=4$, control
channels given by reweighting the prime channels $p\in\{2,3,5,7\}$ and raw flow
$F_0=\partial_{s_0}S$: the spectrum is simple $(0.9619,2.7211,5.6399,8.8009)$,
$\det M=-0.128$ with $\rank M=4$ and condition number $3.4\cdot10^{2}$; the solved
control annihilates all diagonal drifts to $10^{-12}$, and
$\|[S,K]-(F-\Delta(F))\|=1.3\cdot10^{-11}$ with $\|K+K^{T}\|=1.7\cdot10^{-13}$.
\end{remark}

\begin{lemma}[block regulator through collisions]\label{lem:block}
Let $S=S^{*}$ have distinct eigenvalues $\Lambda_\alpha$ with spectral projections
$P_\alpha$ of arbitrary finite multiplicities, and let the flow be controlled as in
Theorem \ref{thm:regulator}, $\dot S=F(u)=F_0+\sum_ru_rF_r$, $F_r=F_r^{*}$. With
\[
  \Delta(F):=\sum_\alpha P_\alpha FP_\alpha,
  \qquad
  K:=\sum_{\alpha\ne\beta}\frac{P_\alpha FP_\beta}{\Lambda_\alpha-\Lambda_\beta},
\]
one has $[S,K]=F-\Delta(F)$ and $K^{*}=-K$. To first order the group $\Lambda_\alpha$
splits by the spectrum of the compression $P_\alpha FP_\alpha$; hence the regulated
condition is the block equation $P_\alpha F(u)P_\alpha=0$ for all $\alpha$, a linear
system in $u$ generalizing $Mu=b$, and when it is solvable the flow $\dot S=[S,K]$ is
isospectral with multiplicities. Between collisions Theorem \ref{thm:regulator}
applies; across a collision the invariant is the inertia (Corollary
\ref{cor:rigidity}); the regulator therefore extends piecewise along any transport
with finitely many collision events.
\end{lemma}

\begin{proof}
$SP_\alpha=\Lambda_\alpha P_\alpha$ and $\sum_\alpha P_\alpha=I$ give
\[
  [S,K]=\sum_{\alpha\ne\beta}(\Lambda_\alpha-\Lambda_\beta)
  \frac{P_\alpha FP_\beta}{\Lambda_\alpha-\Lambda_\beta}
  =F-\Delta(F),
\]
and $K^{*}=-K$ since $F^{*}=F$ and the real denominators are antisymmetric in
$(\alpha,\beta)$. The first-order motion of a degenerate group is degenerate
perturbation theory: the group splits by the eigenvalues of $P_\alpha FP_\alpha$,
which vanish identically iff the block equation holds; the condition is linear in
$u$, and on solution intervals the Lax argument of Theorem \ref{thm:regulator} is
unchanged.
\end{proof}

\section{Warp covariance: the bank-generated readout}\label{sec:warp}

Theorem \ref{thm:regulator} fixes the abstract spectrum $\{\lambda_a\}$ of the regulated
flow; the same transport, unregulated, moves it. The first chart makes reality manifest,
the second carries the physical ordinates, and the map between them --- the warp --- is
not a hypothesis to be assumed but an object the bank constructs. Throughout this
section $t\mapsto(G_0(t),G_1(t))$ is a smooth transport of the continued jet pair, whose
inputs are values and derivatives of the continued logarithmic derivative along the path
--- function data; no zero location enters --- and the parameter interval is
\emph{regular}: $G_0(t)>0$ on the active quotient and $\spec S_t$ simple, where
$S_t:=G_0^{-1/2}G_1G_0^{-1/2}$. The finitely many failures of regularity are the
collision and rank-drop events, crossed piecewise by Lemma \ref{lem:block}.

\begin{definition}[bank velocity field; warp]\label{def:warp}
Let $\lambda_1(t)<\dots<\lambda_n(t)$ be the eigenvalues of $S_t$, realized as the
generalized eigenpairs $G_1w_a=\lambda_aG_0w_a$ normalized by $w_a^{T}G_0w_a=1$, and set
\[
  d_a(t) := w_a^{T}\bigl(\dot G_1-\lambda_a\dot G_0\bigr)w_a
  \qquad\bigl(=\langle v_a,\dot S_t\,v_a\rangle,\ v_a=G_0^{1/2}w_a\bigr).
\]
Let $V_t$ be the unique real polynomial of degree $\le n-1$ with
$V_t(\lambda_a(t))=d_a(t)$, and let the \emph{warp} $\varphi_t$ be the flow of the
scalar equation $\dot x=V_t(x)$, $\varphi_0=\mathrm{id}$.
\end{definition}

\begin{theorem}[warp existence and exact covariance]\label{thm:warp}
On every regular interval the warp is defined on a neighborhood of
$[\lambda_1(0),\lambda_n(0)]$ and satisfies:
\begin{enumerate}
\item[(i)] \textbf{bank determination:} $V_t$, hence $\varphi_t$, is computed from
$(G_0,G_1)(t)$ and $(\dot G_0,\dot G_1)(t)$ alone;
\item[(ii)] \textbf{reality and order:} $\varphi_t(\R\cap\mathrm{dom})\subseteq\R$ and
$x\mapsto\varphi_t(x)$ is strictly increasing, in particular injective;
\item[(iii)] \textbf{exact covariance:} $\varphi_t(\lambda_a(0))=\lambda_a(t)$ for
every $a$ and every $t$ in the interval.
\end{enumerate}
\end{theorem}

\begin{proof}
(i) is the construction. For (ii): $V_t$ has real coefficients --- real nodes, real
diagonal drifts of a Hermitian family --- so the flow preserves $\R$; the field is
polynomial in $x$ with coefficients continuous in $t$, hence locally Lipschitz in $x$
uniformly on compact $t$-intervals, so distinct scalar characteristics never cross and
$\varphi_t$ is strictly increasing where defined. For (iii): first-order perturbation
theory for the Hermitian-definite pencil gives
$\dot\lambda_a=w_a^{T}(\dot G_1-\lambda_a\dot G_0)w_a=d_a(t)=V_t(\lambda_a(t))$, so the
eigenvalue path $t\mapsto\lambda_a(t)$ solves the initial-value problem defining
$\varphi_t(\lambda_a(0))$, and uniqueness makes them equal. Existence: the node
characteristics are the globally defined eigenvalue paths; a point between consecutive
nodes is trapped between two of them by (ii); and continuous dependence on initial
conditions extends the flow to a neighborhood of the compact hull for the same times.
\end{proof}

\begin{theorem}[determinant pullback]\label{thm:pullback}
Let $p_t(x):=\det(S_t-xI)$, and let $I$ be an open real interval containing the
time-$0$ active spectrum on which the warp is defined. Then on $\varphi_t(I)$,
\[
  p_t \;=\; J_t\cdot\bigl(p_0\circ\varphi_t^{-1}\bigr),
\]
with $J_t$ continuous and nonvanishing; consequently
\[
  Z(p_t)\cap\varphi_t(I)\;=\;\varphi_t\bigl(\spec S_0\bigr)\;\subset\;\R .
\]
\end{theorem}

\begin{proof}
The field is polynomial in $x$, so $\varphi_t$ is a strictly increasing $C^1$
diffeomorphism of $I$ onto $\varphi_t(I)$. On $\varphi_t(I)$ the functions $p_t$ and
$p_0\circ\varphi_t^{-1}$ have the same zero set
$\{\varphi_t(\lambda_a(0))\}_a=\{\lambda_a(t)\}_a$ (covariance and injectivity,
Theorem \ref{thm:warp}), each zero simple (simple spectrum on a regular interval). The
ratio $J_t:=p_t/(p_0\circ\varphi_t^{-1})$ is continuous and nonvanishing off the common
zeros and extends across each with the nonzero limit
$p_t'(\lambda_a(t))\big/\bigl(p_0\circ\varphi_t^{-1}\bigr)'(\lambda_a(t))$.
\end{proof}

\begin{remark}[the two architectures, reconciled]\label{rem:architectures}
The regulated and unregulated flows describe one transport in two charts: labels frozen
with a moving readout, or a moving spectrum read through the identity; by Theorem
\ref{thm:warp} the whole discrepancy is $\varphi_t$. This corrects the former statement
of obligation (N), which asked for a spectral equivalence
$\det(G_1-\lambda G_0)=C\det(H^{(1)}_n-\lambda H_n)$ with one spectral variable
$\lambda$ and a constant $C$: under the regulated flow that identity is unsatisfiable
--- the transported spectrum stays at the anchor nodes while the window spectrum sits at
the ordinates, disjoint multisets --- and it is replaced by the pullback identity of
Theorem \ref{thm:pullback}, in which the change of spectral variable is explicit and
bank-generated.
\end{remark}

\begin{proposition}[the warp is internally generated on the anchor segment]\label{prop:selfcal}
On the segment $s_0>1$ the transport is $\dot\sigma=-x\,\sigma$ (each prime clock damps
by its own frequency), so $\dot m_k=-m_{k+1}$: the moment flow is the pencil's own
shift. Let $(\lambda_a,w_a)$ be the Gauss data of the $n$-point pencil, $P_n$ the monic
node polynomial, and $h_n:=\int P_n^2\,d\sigma=m_{2n}-\sum_aw_a\lambda_a^{2n}$ the
quadrature error of $x^{2n}$. Then
\[
  \frac{d\lambda_a}{ds_0}\;=\;-\,\frac{h_n}{w_a\,P_n'(\lambda_a)^2}\,,
\]
strictly negative for every $a$. Every quantity on the right except $h_n$ is internal
to the pencil; the entire external input of the flow is the single boundary scalar
$h_n$.
\end{proposition}

\begin{proof}
The $n$-point Gauss data matches the moments $m_k$ for $k\le2n-1$. Differentiating the
matching conditions along $\dot m_k=-m_{k+1}$, the homogeneous solution
$\dot w_a=-\lambda_aw_a$, $\dot\lambda_a=0$ accounts for $-m_{k+1}$ within quadrature
exactness, leaving the forcing $-h_n\delta_{k,2n-1}$ at the single degree where
exactness fails. The correction $(u_a,\dot\lambda_a)$ solves
$\sum_a[u_aQ(\lambda_a)+w_a\dot\lambda_aQ'(\lambda_a)]=-h_n\cdot[x^{2n-1}]Q$ for all
$Q$ of degree $\le2n-1$. Take $Q=\ell_aP_n$ with $\ell_a$ the Lagrange basis: $Q$
vanishes at every node, $Q'(\lambda_b)=\delta_{ab}P_n'(\lambda_a)$, and
$[x^{2n-1}]Q=1/P_n'(\lambda_a)$; the display follows. Positivity of $\sigma$ on the
segment gives $h_n>0$, $w_a>0$.
\end{proof}

\begin{remark}[self-calibration, and its limit]\label{rem:selfcal}
Proposition \ref{prop:selfcal} is the Toda-deformation structure of exponentially
damped measures: the infinite operator is isospectral --- its spectrum is the fixed
prime-clock support $\{\log n\}$, the carrier --- while the finite pencil's nodes, the
readout, move only through the boundary forcing $h_n$. Verified
(\texttt{tmp/att205\_selfcal\_probe.py}, $s_0=1.5$, $N=4$): the internal law matches
the direct perturbation drifts to $1.6\cdot10^{-11}$, with $h_n=2.2\cdot10^{5}$
cancelling against Christoffel factors of order $10^{-4}$ exactly. Two consequences
and one boundary. First, the warp needs no external support on the anchor segment
beyond one scalar per instant; the continuation's defects (pole, $\Gamma$, window)
must enter the finite system through this single boundary channel, so Hypothesis
\ref{hyp:seat} becomes a statement about the behavior of a one-dimensional boundary
sequence rather than an $n\times n$ identity. Second, $h_n=\int P_n^2\,d\sigma$ is
manifestly nonnegative exactly as long as the driving functional is a positive
measure: where the continuation costs $h_n$ its manifest sign is where the seat's
content begins. The boundary: a self-calibrating flow is not a self-certifying one
--- the device generates its own motion, but the terminal seating certificate must
still come from the crossing analysis of Lemmas \ref{lem:nullflow} and
\ref{lem:resonance}. Finally, an identity rather than an analogy: for Hankel pairs,
$h_n=\det H_{n+1}/\det H_n$, so the boundary scalar that drives the flow \emph{is}
the inertia-increment detector of the census --- its loss of sign at one size and
the appearance of a negative direction at the next are the same fact (verified along
the continuation: sign flip at $n=14$, inertia event at $n=15$;
\texttt{tmp/att211\_boundary\_onset.py}).
\end{remark}

\begin{remark}[numerical verification]
On the convergent segment ($s_0\colon1.5\to1.2$, $N=4$, exact jets of $-\zeta'/\zeta$;
\texttt{tmp/att201\_warp\_probe.py}): the perturbation drift matches finite differences
of the eigenvalue paths to $8.9\cdot10^{-8}$; the flow of the interpolated field carries
the anchor nodes $(1.1916,\,4.0395,\,9.6214,\,19.3384)$ onto the terminal eigenvalues
$(2.1714,\,9.2832,\,23.2371,\,47.5293)$ with maximal error $1.9\cdot10^{-10}$; an
interior test point remains strictly interior. The exact-jet anchor nodes differ from
those quoted after Theorem \ref{thm:regulator}, which realize the truncated prime sum at
finite cutoff; the two banks are distinct legitimate realizations, and the drift between
them is the truncation observation of the ledger.
\end{remark}

\section{The conditional theorem}

\begin{hypothesis}[$\mathrm{PSD}$]\label{hyp:psd}
For every window $W$ of a tiling of $\R$ and $n=\lceil\mu_0(W)\rceil+1$,
$H_n(W)\succeq0$; equivalently (Proposition \ref{prop:stieltjes}) both Hankel matrices of
the even moments are positive semidefinite in the quotient chart.
\end{hypothesis}

\begin{theorem}\label{thm:main}
Assume Hypothesis \ref{hyp:psd}. Then
\begin{enumerate}
\item every nontrivial zero of $\zeta$ lies on the critical line;
\item for each window the spectral data of $H_n(W)$ returns the full census of Theorem
\ref{thm:census}, with $n_-\equiv0$;
\item the zeros are the real generalized eigenvalues of the explicitly computed
Hermitian pencils $\bigl(H^{(1)}_n(W),H_n(W)\bigr)$.
\end{enumerate}
\end{theorem}

\begin{proof}
By Theorem \ref{thm:inertia}, $H_n(W)\succeq0$ with $n\ge m(W)$ forces $q(W)=0$: no
conjugate pair of zeros of $A$ lies in $W$. Nonreal zeros occur in conjugate pairs with
equal real part, hence both members lie in the same window of the tiling (for a pair on a
boundary, apply the argument to a shifted tiling). Every nontrivial zero of $\zeta$
corresponds to a zero of $A$ lying in some window, and only finitely many per window by
Riemann--von Mangoldt. Hence all zeros of $A$ are real, i.e.\ all nontrivial zeros of
$\zeta$ have real part $\tfrac12$. Items (2) and (3) are Theorem \ref{thm:census}, whose
pencil is Hermitian-definite on $\ran H_n$ under Hypothesis \ref{hyp:psd}.
\end{proof}

\section{The seat: reduction to a single hypothesis}\label{sec:seat}

\begin{definition}[terminal defect]\label{def:defect}
Fix a window $W$ with matched dimension $n$ and a transport path whose terminal chart is
the window chart, and let $(G_0(1),G_1(1))$ be the terminal transported pair. Put
\[
  D_0 := G_0(1)-H_n(W),\qquad D_1 := G_1(1)-H^{(1)}_n(W).
\]
Both differences are computed from function data --- jets of the continued logarithmic
derivative on one side, contour integrals of $A'/A$ on the other --- so the defect is
defined with no reference to zeros. Obligation (T) of \S\ref{sec:obligations} evaluates
$D_\ell$ in closed form; the smooth blocks are now explicit. The pole block is the
Hankel of the exponential measure $e^{-(s_0-1)x}dx$ on $[0,\infty)$, moments
$k!/(s_0-1)^{k+1}$, positive semidefinite on the convergent segment $s_0>1$. The
$\Gamma$-block (trivial zeros) is the Hankel of the measure
$e^{-(s_0+2)x}\bigl(1-e^{-2x}\bigr)^{-1}dx$ on $(0,\infty)$, with moments in closed
Hurwitz form
\[
  k!\;2^{-(k+1)}\,\zeta\!\left(k+1,\;1+\tfrac{s_0}2\right)\qquad(k\ge1),
\]
a manifestly positive density for every $s_0>-2$, hence positive semidefinite
\emph{throughout the strip}. The entire smooth part of the defect therefore adds
positivity at every station of the transport; only the window and lateral blocks of
(T) remain unevaluated. (These closed forms are the trivial-zero and pole components
of the jet engine used in the numerical maps, validated there against direct
differentiation to $5\cdot10^{-5}$.)
\end{definition}

\begin{lemma}[null-cone flow decomposition]\label{lem:nullflow}
Let $\sigma=\sum_{x\in X}w(x)\delta_x$ be an atomic functional on a conjugation-closed
multiset $X=\overline X$ with weights satisfying $w(\bar x)=\overline{w(x)}$, and let
$H$ be the Hankel matrix of its moments, of size $n$. For every real polynomial $P$ of
degree $<n$ with coefficient vector $v$,
\[
  v^{T}Hv \;=\; \sum_{x\in X\cap\R}w(x)P(x)^{2}
  \;+\!\!\sum_{\{z,\bar z\}\subset X\setminus\R}\!\!2\,\mathrm{Re}\bigl[w(z)P(z)^{2}\bigr],
\]
in which the real-support sum is termwise nonnegative whenever $w\ge0$ on $X\cap\R$,
while the conjugate-pair terms are unrestricted in sign. Consequently, along a weight
homotopy $t\mapsto w_t$ with $\dot w_t\ge0$ on $X\cap\R$, the flow through a null
direction ($v^{T}H_tv=0$) obeys
\[
  v^{T}\dot H_tv \;\ge\; \sum_{\{z,\bar z\}}2\,\mathrm{Re}\bigl[\dot w_t(z)P(z)^{2}\bigr]:
\]
a negative-sector crossing must be fed by conjugate-pair terms; real support cannot
produce one. In particular, on any pointwise-monotone kernel homotopy whose weighted
support is real and which starts at a nonnegative-weight configuration,
$n_-(H_t)\equiv0$ unconditionally.
\end{lemma}

\begin{proof}
Group the atoms: conjugate pairs combine to $2\mathrm{Re}[w(z)P(z)^{2}]$ because $P$
is real, and the real-support terms are $w(x)P(x)^2\ge0$ termwise. The flow statement
is the same identity applied to $\dot H_t$, discarding the nonnegative real sum. The
last claim is closedness of the PSD cone: positive start, and no crossing is
available.
\end{proof}

\begin{remark}[measured margin law]
On the verified window $W=(10,32)$ ($n=5$, four zeros;
\texttt{tmp/att202\_margin\_probe.py}), along the error-function smoothing of the
window indicator: the dead-direction mass $v_*^{T}H(\tau)v_*$ ($v_*$ the terminal null
polynomial vanishing at the four interior zeros) is monotone with positive flow sign
($+1.4\cdot10^{-2}$ at $\tau=1$); by $\tau=1$ its source is $100\%$ the
boundary-nearest exterior zero $\gamma_5=32.935$; and $\lambda_{\min}(H(\tau))$ tracks
the Rayleigh value $v_*^{T}Hv_*$ to $99.8\%$ at $\tau=0.25$. The positivity slack of
the window pencil near the sharp limit is the weight leaking onto the nearest exterior
zero through the terminal null polynomial: the margin has a single computable source.
\end{remark}

\begin{remark}[cosine transport of the pair-fed flow]\label{rem:cosine}
The conjugate-pair terms of Lemma \ref{lem:nullflow} have an exact operator form: for
$F$ entire and real on $\R$, $F(x+iy)+F(x-iy)=2\cos(y\partial_x)F(x)$ (Taylor
expansion, absolutely convergent). Hence a pair at depth $y$ feeds the flow with the
\emph{same} real density $g=\dot w\,P^2$ that a real zero would, transported by
$\cos(y\partial_x)$. The Fourier multiplier $\cos(y\xi)$ is nonnegative for
$|\xi|\le\pi/(2y)$, and $y<\tfrac12$ unconditionally (the critical strip), so the
worst-case positivity radius of the transport is the band $|\xi|\le\pi$. Two
consequences. First, a quantitative target: pair-fed dominance on a null direction
requires the density $g$ to carry Fourier mass beyond the positivity radius of the
actual pair depths present --- shallow pairs feed like real zeros. Second --- correcting an identification staked in an earlier draft of this remark ---
the cross-check against the narrow-support literature \emph{fails on the numbers, and
informatively so}: the classical unconditional positivity radius (Yoshida
\cite{Yoshida1992}) is
$\log2=0.693$, set by the \emph{prime vacuum} (no $\Lambda$-support below the first
prime), whereas the cosine-transport radius is $\pi=3.14$, set by pair depth. The two
are distinct mechanisms on opposite sides of the explicit formula: prime-vacuum limits
the classical regime, pair depth limits this transport, and $\pi>\log2$ means the
pair-depth constraint is strictly the weaker one --- positivity of the pair-fed flow
persists beyond the prime-vacuum boundary, and the two constraints compose rather
than coincide. The present formulation additionally needs the sign only on null
directions (a thinner set) and may weigh actual pair depths against unconditional
density estimates rather than assuming the worst case everywhere.
\end{remark}

\begin{lemma}[resonance necessity]\label{lem:resonance}
Let $P$ be a real polynomial of degree $d$ with roots $r_1,\dots,r_d$ (nonreal ones in
conjugate pairs), let $z_0=x_0+iy_0$ with $y_0>0$, and let $w(z_0)\in\C^{\times}$ with
$|\arg w(z_0)|\le\varphi_w<\tfrac\pi2$. Define the deviation of each root as the angle
between $z_0-r_j$ and the real line,
\[
  \delta_j \;:=\; \arctan\!\frac{|y_0-\mathrm{Im}\,r_j|}{|x_0-\mathrm{Re}\,r_j|}
  \;\in\;(0,\tfrac\pi2]
  \qquad(\delta_j:=\tfrac\pi2\ \text{if}\ x_0=\mathrm{Re}\,r_j).
\]
If $\mathrm{Re}\bigl[w(z_0)P(z_0)^2\bigr]\le0$, then
\[
  \sum_{j=1}^{d}\delta_j\;\ge\;\frac\pi4-\frac{\varphi_w}2,
\]
and in particular some root satisfies
\[
  |x_0-\mathrm{Re}\,r_j|\;\le\;
  \frac{|y_0-\mathrm{Im}\,r_j|}{\tan\bigl((\tfrac\pi4-\tfrac{\varphi_w}2)/d\bigr)}
  \;\approx\;\frac{4d}{\pi}\,\bigl(y_0+|\mathrm{Im}\,r_j|\bigr)
  \quad\text{for small angles.}
\]
\end{lemma}

\begin{proof}
Each squared factor of $P(z_0)^2=c^2\prod_j(z_0-r_j)^2$ contributes, modulo $2\pi$, an
argument $\pm2\delta_j$: for a factor with argument $\theta_j\in(0,\pi)$, either
$\theta_j=\delta_j$ or $\theta_j=\pi-\delta_j$, and $2(\pi-\delta_j)\equiv-2\delta_j$.
Hence $\arg\bigl[P(z_0)^2\bigr]\equiv\sum_j\pm2\delta_j \pmod{2\pi}$, and
$\bigl|\sum_j\pm2\delta_j\bigr|\le2\sum_j\delta_j$. A nonpositive real part of
$w(z_0)P(z_0)^2$ requires its total argument to lie within $\tfrac\pi2$ of $\pi$
modulo $2\pi$, so $2\sum\delta_j+\varphi_w\ge\tfrac\pi2$. The root bound is the
pigeonhole $\max_j\delta_j\ge(\tfrac\pi4-\tfrac{\varphi_w}2)/d$.
\end{proof}

\begin{corollary}[separation forces depth]\label{cor:separation}
In the setting of Lemma \ref{lem:resonance}, suppose every root satisfies
$|x_0-\mathrm{Re}\,r_j|\ge d$ and $|\mathrm{Im}\,r_j|\le\iota$. Then a nonpositive
pair term at $(x_0,y_0)$ requires
\[
  y_0\;\ge\;d\,\tan\!\frac{\tfrac\pi4-\tfrac{\varphi_w}2}{\,d^{\circ}\!P\,}\;-\;\iota
  \;\approx\;\frac{\pi d}{4\,d^{\circ}\!P}\quad(\varphi_w,\iota\ \text{small}),
\]
where $d^{\circ}\!P$ is the degree. Equivalently: the abscissas from which a pair of
depth $y_0$ can cross are confined to the
$\bigl((y_0+\iota)/\tan((\tfrac\pi4-\tfrac{\varphi_w}2)/d^{\circ}\!P)\bigr)$-neighborhood
of the null-root abscissa set --- a set of total measure at most
$\approx\tfrac{8}{\pi}\,(d^{\circ}\!P)^2\,y_0$ --- while pairs elsewhere in the window
are spectators regardless of the flow.
\end{corollary}

\begin{proof}
Each deviation obeys $\delta_j\le\arctan((y_0+\iota)/d)$, so
$\sum_j\delta_j\le d^{\circ}\!P\cdot\arctan((y_0+\iota)/d)$; Lemma \ref{lem:resonance}
requires this to reach $\tfrac\pi4-\tfrac{\varphi_w}2$. The neighborhood statement is
the same inequality read per root, and the measure bound is $d^{\circ}\!P$ intervals
of the stated half-width.
\end{proof}

\begin{remark}[the composition slot]
Corollary \ref{cor:separation} is the deterministic half of a two-part exclusion:
crossing-capable pairs at depth $y$ are confined to an explicit set of measure
$O(n^2y)$ per window, while the number of pairs of depth $\ge y$ up to height $T$ is
counted unconditionally by the zero-density estimates: $N(\sigma,T)\ll
T^{3(1-\sigma)/(2-\sigma)+o(1)}$ (Ingham \cite{Ingham1940}) and
$N(\sigma,T)\le T^{30(1-\sigma)/13+o(1)}$ (Guth--Maynard \cite{GuthMaynard2024}, the
first exponent improvement in eighty-four years), so pairs at depth $y$ number
$\ll T^{30(\frac12-y)/13+o(1)}$. Two calibrations. The composition of the two counts
against the window census yields rarity statements --- windows admitting a resonance
are sparse --- and rarity is not exclusion: a count cannot reach zero, so this route's
unaided yield is a new-mechanism density theorem, not the seat. Closing the seat on a
given window requires the per-window exclusion --- path steering within the vulnerable
set of Corollary \ref{cor:separation}, or positivity --- and that is where the
remaining content sits.
\end{remark}

\begin{remark}[crossings are local resonances]
Combining Lemmas \ref{lem:nullflow} and \ref{lem:resonance}: a negative-sector
crossing requires a conjugate pair at depth $y_0<\tfrac12$ whose abscissa lies within
$\approx\frac{4(n-1)}{\pi}\,y_0$ of a root cluster of the null polynomial of the
accumulated Hankel --- with the kernel's phase budget $\varphi_w$ a design parameter
one makes small on the strip. Shallow pairs can cross only in near-collision with a
null root (the continuity of Corollary \ref{cor:rigidity}, now quantitative); deep
pairs have room, but their number is the subject of unconditional zero-density
estimates. The seat's remaining content is the exclusion of these localized
resonances along one designed path per window; the composition of depth, density and
the locality of null directions (the measured margin law) is the open program.
\end{remark}

\begin{hypothesis}[the seat $(\mathrm{S})$]\label{hyp:seat}
For every window $W$ of the tiling there is a transport path with finitely many
regularity events such that every point of
$\spec\bigl(H^{(1)}_n(W),H_n(W)\bigr)\big|_{\ran H_n}$ --- the window support --- is a
terminal limit of the warped anchor spectrum $\varphi_t(\spec S_0)$, the remaining
$n-m$ directions retiring into the null space through the rank-drop events.
\end{hypothesis}

\begin{remark}[three contractions of one kernel]\label{rem:contractions}
The transported objects of this paper are contractions of the single two-variable
kernel $B_{jk}(\rho,\rho')=(s_0-\rho)^{-(j+1)}(s_0-\rho')^{-(k+1)}$ along three
pairings of the conjugation- and FE-closed zero multiset: the \emph{holomorphic
diagonal} ($\rho$ with $\rho$; the jet-Hankel face, real symmetric), the
\emph{functional-equation pairing} ($\rho$ with $1-\rho$; hermitian by
conjugation-closure), and the \emph{conjugate pairing} ($\rho$ with $\bar\rho$; a
Gram matrix, positive semidefinite always). Atomwise, the FE and conjugate pairings
coincide iff $\bar\rho=1-\rho$, i.e.\ iff the zero lies on the critical line; the
holomorphic diagonal is indefinite for complex frequencies whether or not they lie on
the line. Two consequences. First, a calibration: the indefiniteness onset of the
jet-Hankel face along the continuation carries no information about zero locations
--- it is the generic behavior of complex-frequency moment data. Second, the
diagnostic content of the program is exactly the coincidence of the FE and conjugate
contractions --- every zero self-paired by the functional equation is a zero on the
line --- which is the pairing-language form of the single hypothesis, and the
positivity of the FE contraction is its window-free shadow.
\end{remark}

\begin{proposition}[scalar seat criterion]\label{prop:scalarseat}
Over the nontrivial-zero multiset put
\[
  S(s)\;:=\;\sum_{\rho}\,(s-\rho)^{-1}\,(\bar s-1+\rho)^{-1},
\]
absolutely convergent for $s$ off the zero set (terms are $O(|\gamma_\rho|^{-2})$),
and real-valued: the relabeling $\rho\mapsto1-\bar\rho$, under which the multiset is
closed, conjugates its terms. Then
\[
  \text{every nontrivial zero lies on the critical line}
  \iff
  S(s)\ge0\ \text{ for all }s\ \text{in the strip off the zero set.}
\]
\end{proposition}

\begin{proof}
($\Rightarrow$) For $\rho$ on the line, $1-\rho=\bar\rho$, so the $\rho$-term is
$(s-\rho)^{-1}\overline{(s-\rho)^{-1}}=|s-\rho|^{-2}>0$: the sum is termwise
positive. ($\Leftarrow$) Let $\rho_0=\beta_0+i\gamma_0$ with $\beta_0\ne\tfrac12$ and
set $s=\rho_0+\varepsilon e^{i\theta}$. The $\rho_0$-term is
$\varepsilon^{-1}e^{-i\theta}\bigl((2\beta_0-1)+\varepsilon e^{-i\theta}\bigr)^{-1}$,
of modulus $\sim(\varepsilon|2\beta_0-1|)^{-1}$ with freely rotating phase; choosing
$\theta$ with $\cos\theta\cdot(2\beta_0-1)<0$ makes its real part $\to-\infty$ as
$\varepsilon\to0$. Every other term stays bounded on the $\varepsilon$-ball: the
remaining members of the $\rho_0$-family ($\bar\rho_0$, $1-\rho_0$, $1-\bar\rho_0$)
and all other zeros are at distance bounded below, and the full sum of their moduli
converges. Hence $S(s)<0$ for small $\varepsilon$.
\end{proof}

\begin{remark}[the traveling seat]\label{rem:travelingseat}
Proposition \ref{prop:scalarseat} is the $n=1$ member of the FE-paired family
$A_{jk}(s)=\sum_\rho(s-\rho)^{-(j+1)}(\bar s-1+\rho)^{-(k+1)}$, Hermitian for every
anchor $s$ by the same relabeling, with on-line zeros contributing rank-one positive
terms at every anchor. Numerically (\texttt{tmp/att217\_anchor\_lift.py}, verified
range, $n=8$): as the anchor rises through the strip the form remains positive
semidefinite at every station and \emph{seats} on each zero in turn --- mass share of
the nearest zero $99\%$ at $\tau=\gamma_j$, resolvent spike at each seating,
conditioning improving with localization. A negative direction would appear exactly
at the seating of an off-line pair: the falsifier is localized in the anchor height.

The source reading required by the working rules identifies the scalar criterion: by
partial fractions, $(s-\rho)^{-1}(\bar s-1+\rho)^{-1}
=(s+\bar s-1)^{-1}\bigl[(s-\rho)^{-1}+(\rho-(1-\bar s))^{-1}\bigr]$, and summing over
the paired multiset via the Hadamard expansion and the functional equation
($\xi'/\xi(1-\bar s)=-\overline{\xi'/\xi(s)}$) gives the closed form
\[
  S(s)\;=\;\frac{2\,\mathrm{Re}\,\bigl[\xi'/\xi(s)\bigr]}{2\sigma-1},
  \qquad \sigma=\mathrm{Re}\,s .
\]
Proposition \ref{prop:scalarseat} is therefore equivalent to the positivity
criterion of Hinkkanen (1997) and Lagarias \cite{Lagarias1999}: RH iff
$\mathrm{Re}\,[\xi'/\xi](s)>0$ for $\sigma>\tfrac12$ --- a rediscovery, in
FE-paired coordinates, with the pairing trichotomy of Remark \ref{rem:contractions}
supplying a two-line derivation. What is retained beyond the identification: the
anchor-lift geometry above (the criterion's falsifier is seat-localized), and the
wiring of the criterion into the transport apparatus of this paper. The graded
family $A_{jk}$ for $n>1$ is likewise identified by the required source reading: it
is of the Weil--Yoshida hermitian-form species \cite{Yoshida1992,Bombieri2000} in a
resolvent basis at complex anchor --- the species carrying Li's criterion as its
Laguerre-basis instance (Bombieri--Lagarias \cite{BombieriLagarias1999}) and the
Connes--Consani semi-local positivity results as its deepest partial positivity.
The bridge runs both ways: partial positivity results for these forms are partial
results on the seat, and the transported-pencil apparatus --- the internal flow law,
the boundary-scalar detector, the seat-localized falsifier geometry --- is what this
paper adds around the classical species.
\end{remark}

\begin{proposition}[unconditional positivity on the verified band]\label{prop:verifiedband}
Let $T_0$ be a height below which every nontrivial zero lies on the critical line
($T_0=3\cdot10^{12}$ by \cite{PlattTrudgian2021}). There is an absolute constant $C$
such that
\[
  S(s)\;>\;0
  \qquad\text{for every }s=\sigma+i\tau\text{ in the strip with }
  30\le|\tau|\le T_0-C\log T_0 .
\]
The scalar seat criterion of Proposition \ref{prop:scalarseat} therefore holds
unconditionally on the full anchor band of the verified range; its open content is
confined to anchors above the verification height.
\end{proposition}

\begin{proof}
Split the sum at height $T_0$. Below: every zero is on the line, so each term is
$|s-\rho|^{-2}>0$; by Riemann--von Mangoldt, $[\,|\tau|,|\tau|+1\,]$ contains a zero
once $|\tau|\ge30$, giving a single-term lower bound
$S_{<}(s)\ge\bigl(\tfrac14+1\bigr)^{-1}\cdot\tfrac12$, say. Above: for any zero
$\rho$ (on the line or not) at height $\gamma>T_0$, both factors have modulus at
least $\gamma-|\tau|$ --- $|s-\rho|\ge\gamma-|\tau|$ and
$|\bar s-1+\rho|\ge\gamma-|\tau|$, since each has imaginary part $\pm(\gamma-|\tau|)$
up to sign conventions --- so the term is bounded by $2(\gamma-|\tau|)^{-2}$ in
modulus. Integrating against the counting measure,
\[
  |S_{>}(s)|\;\le\;2\int_{T_0}^{\infty}\frac{dN(\gamma)}{(\gamma-|\tau|)^{2}}
  \;\ll\;\frac{\log T_0}{\,T_0-|\tau|\,},
\]
which is smaller than the lower bound once $T_0-|\tau|\ge C\log T_0$ for an absolute
$C$.
\end{proof}

\begin{proposition}[band-local seat criterion]\label{prop:bandlocal}
For $H\ge60$, with $C$ the absolute constant of Proposition \ref{prop:verifiedband}:
\begin{enumerate}
\item[(i)] if every zero with $|\gamma|\le H$ lies on the line, then $S(s)>0$ for all
anchors with $30\le|\tau|\le H-C\log H$ (the proof of Proposition
\ref{prop:verifiedband} verbatim, with $T_0$ replaced by $H$);
\item[(ii)] if $S\ge0$ on the band $|\tau|\le H$, then every zero with
$|\gamma|\le H-1$ lies on the line (the converse construction of Proposition
\ref{prop:scalarseat} produces a negative value at an anchor within distance $1$ of
any off-line zero).
\end{enumerate}
In particular RH is equivalent to band positivity at every height, with the
correspondence exact up to logarithmic boundary layers.
\end{proposition}

\begin{remark}[the induction frame]\label{rem:induction}
Proposition \ref{prop:bandlocal} places the whole question in one-dimensional
coordinates: RH holds iff the band-positivity property is unbounded, and by (U4) it
holds below $3\cdot10^{12}$. The single remaining statement is the self-extension
step --- band positivity below $H$ implies it below $H+C\log H$ --- and by (ii) this
step is equivalent to the absence of a first off-line zero in the boundary layer:
necessarily of full strength, as every reduction in this paper must be, but now
localized to a moving layer of logarithmic width, where the protecting mass is the
verified band below and the threat is a single hypothetical seating just above. The
seat, in its minimal coordinates: \emph{the boundary layer never turns.}
\end{remark}

\begin{theorem}[the chain]\label{thm:chain}
Hypothesis \ref{hyp:seat} implies Hypothesis \ref{hyp:psd}. Hence, granted the seat
alone: every nontrivial zero of $\zeta$ lies on the critical line, the census of
Theorem \ref{thm:census} is returned on every window, and the zeros are the real
spectra of the computed Hermitian pencils (Theorem \ref{thm:main}).
\end{theorem}

\begin{proof}
On every regular interval the transported spectrum is real: the anchor spectrum is real
(Proposition \ref{prop:anchor}, $S_0$ Hermitian) and Theorem \ref{thm:warp} transports
it by a real map; across regularity events reality is preserved because the pencil
remains Hermitian with $G_0>0$ on the active quotient (Lemma \ref{lem:block}). A limit
of reals is real, so by Hypothesis \ref{hyp:seat} every support point of the window
measure \eqref{eq:moments} is real: $q(W)=0$ for every window, and Theorem
\ref{thm:inertia} gives $H_n(W)\succeq0$: Hypothesis \ref{hyp:psd}.
\end{proof}

\begin{remark}[necessity, and what has been gained]
By the rigidity of the moment problem (\S\ref{sec:obligations}) any hypothesis implying
Hypothesis \ref{hyp:psd} is equivalent to it; the seat is no exception, and this is
forced --- it is not a defect of the method. What has changed is the form: of the former
obligations, (N), (W1), (W2) and (C) are now theorems (Theorem \ref{thm:regulator},
Lemma \ref{lem:block}, Theorems \ref{thm:warp} and \ref{thm:pullback}), and the residue
is one statement, bank-generated and zero-free in formulation, finite-dimensional on
each window.
\end{remark}

\section{Remaining obligations}\label{sec:obligations}

Theorem \ref{thm:main} is conditional on Hypothesis \ref{hyp:psd}, and Theorem
\ref{thm:chain} reduces Hypothesis \ref{hyp:psd} to the seat. By the rigidity of the
moment problem, any pairing whose Gram matrix reproduces the moments $\mu_{j+k}(W)$
\emph{is} the counting pairing of \eqref{eq:moments}, so its positivity is equivalent to
the conclusion; a proof must therefore proceed from the bank side, and the precise
remaining steps are these.

\begin{enumerate}
\item[(T)] \textbf{Contour transport (defect evaluation).} Evaluate the terminal defect
of Definition \ref{def:defect} in closed form,
\[
  G_\ell[j,k] \;=\; \frac{1}{2\pi i}\oint_\Gamma z^{\,j+k+\ell}\frac{A'(z)}{A(z)}\,dz
  \;+\; D_\ell[j,k],
\]
with $D_\ell$ the explicit pole, $\Gamma$-factor, lateral and windowing contributions of
(U5). The chain of Theorem \ref{thm:chain} does not consume (T): the defect is already
defined as a difference of computed objects; (T) supplies its closed form, the DC
structure of the seat, and computational access to Hypothesis \ref{hyp:seat} on explicit
windows.
\item[(S)] \textbf{The seat.} Hypothesis \ref{hyp:seat}: terminal seating of the warped
anchor spectrum on the window support. This is the single remaining hypothesis. By
Lemma \ref{lem:nullflow} its content along monotone kernel homotopies is exactly: the
conjugate-pair-fed component of the flow never dominates the real-fed component on a
null direction of the accumulated Hankel --- the only crossing mechanism that exists. The
former obligations (N), (W) and (C) are discharged: pencil neutrality by Theorem
\ref{thm:regulator} with the corrected terminal statement of Remark
\ref{rem:architectures}; warp covariance --- bank determination, reality, injectivity,
and the determinant-pullback law --- by Theorems \ref{thm:warp} and \ref{thm:pullback};
collisions by Lemma \ref{lem:block}.
\end{enumerate}

Granted (S), the chain runs forward: the Euler-anchored pencil is positive semidefinite
unconditionally (Proposition \ref{prop:anchor}); its spectrum is real and is carried by
the bank-generated warp along the transport, exactly and order-preservingly (Theorems
\ref{thm:regulator}, \ref{thm:warp}, \ref{thm:pullback}, Lemma \ref{lem:block}); the
seat identifies the terminal limits with the window support; reality forces $q\equiv0$
(Theorem \ref{thm:inertia}); and Theorem \ref{thm:main} concludes.

\begin{thebibliography}{9}

\bibitem{Ingham1940} A.~E. Ingham, \emph{On the estimation of $N(\sigma,T)$},
Quart. J. Math. Oxford 11 (1940), 291--292.

\bibitem{GuthMaynard2024} L.~Guth and J.~Maynard, \emph{New large value estimates
for Dirichlet polynomials}, arXiv:2405.20552 (2024).

\bibitem{PlattTrudgian2021} D.~Platt and T.~Trudgian, \emph{The Riemann hypothesis
is true up to $3\cdot10^{12}$}, Bull. London Math. Soc. 53 (2021), 792--797.

\bibitem{Yoshida1992} H.~Yoshida, \emph{On hermitian forms attached to zeta
functions}, in: Zeta Functions in Geometry, Adv. Stud. Pure Math. 21 (1992).

\bibitem{Lagarias1999} J.~C. Lagarias, \emph{On a positivity property of the
Riemann $\xi$-function}, Acta Arith. 89 (1999), 217--234.

\bibitem{BombieriLagarias1999} E.~Bombieri and J.~C. Lagarias, \emph{Complements
to Li's criterion for the Riemann hypothesis}, J. Number Theory 77 (1999),
274--287.

\bibitem{Bombieri2000} E.~Bombieri, \emph{Remarks on Weil's quadratic functional
in the theory of prime numbers, I}, Rend. Mat. Acc. Lincei (9) 11 (2000),
183--233.

\end{thebibliography}

\end{document}
